\documentclass[12pt,oneside,a4paper]{article}
\usepackage{amssymb} % blackbaord bold 
\usepackage{amsmath} % for modulo function
\usepackage{hyperref}
\usepackage{xcolor}
\usepackage{listings}
\usepackage{graphicx}
\usepackage{csvsimple}
\usepackage{placeins}
\usepackage{array}
\usepackage{booktabs}
\author{Qianqian Shan}
\title{SQL Notes}
\date{\vspace{-5ex}}
\begin{document}
\maketitle
\tableofcontents
\newpage


\lstdefinestyle{customc}{
	belowcaptionskip=1\baselineskip,
	breaklines=true,
	%frame=shadowbox,
	xleftmargin=\parindent,
	language=C,
	showstringspaces=false,
	basicstyle=\footnotesize\ttfamily,
	keywordstyle=\bfseries\color{green!40!black},
	commentstyle=\itshape\color{purple!40!black},
	identifierstyle=\color{black},
	stringstyle=\color{orange},
}

\lstset{escapechar=none,style=customc}

\lstset{language=SQL}          % Set your language (you can change the language for each code-block optionally)
\section{Basic Concepts}
\textbf{Database}: A container (usually a file or a set of files) to store organized data. 
\\~\\
\textbf{Table}: A structured list of data of a specific type. Every table in the database has a \emph{unique} name that identifies it. 
\\~\\
\textbf{Schema}: Information about database and table layout and properties. For example, information that how data is stored in them, what data is stored , how it's broken up etc. 
\\~\\
\textbf{Primary Keys}: A column(or a set of columns) whose value \emph{uniquely} identify every row in a table. \\~\\
\textbf{Wildcards}: Special characters used to match parts of a value. For example, ``\%'', ``\_''.\\~\\
\textbf{Foreign Key}: A column in one table that contains the primary key from another table, thus defining the relationships between tables. 

\section{Retrieving Data}
\subsection{Display information}
\begin{lstlisting}
-- Return a list of available databases.
SHOW DATABASES; 

-- Show a list of tables within a database 
SHOW TABLES;

-- Show columns from a table 
SHOW COLUMNS FROM customers;

-- Use describe as a shortcut of show...from...
DESCRIBE customers;

\end{lstlisting}

\subsection{The \textsf{SELECT} statement} 
% \texttt{SELECT}

\begin{lstlisting}

SELECT prod_name,prod_id FROM products;

-- Select all columns 
SELECT * FROM products;

-- Retrieve distinct rows 
SELECT DISTINCT vend_id FROM products; 
--DISTINCT applies to all columns, not just the one it precedes.

-- Limit the results by LIMIT 
SELECT prod_name FROM products LIMIT 5;


-- Return five rows starting from row 4.
SELECT prod_name FROM products LIMIT 4,5;

-- Or alternatively, 
SELECT prod_name FROM products LIMIT 5 OFFSET 4;
\end{lstlisting}

\subsection{Sort Retrieved Data}
1. Sort data by \textsf{ORDER BY}. \\
\begin{lstlisting}
SELECT prod_name 
FROM products 
ORDER BY prod_name;

-- Sort by multiple columns 
SELECT prod_id, prod_price 
FROM products 
ORDER BY prod_price DESC, prod_id;
\end{lstlisting} 

%\setlength{\parindent}{0.3em}
\noindent  % surpress paragraph indentation  
2. Specify sort direction by \textsf{DESC}, \textsf{ASC} is the default order. \textsf{DESC} applies to the column which directly precedes it. 
\begin{lstlisting}
SELECT prod_id, prod_price, prod_name 
FROM products 
ORDER BY prod_price DESC;

-- Sort by multiple columns.
SELECT prod_id, prod_price, prod_name
 FROM products 
ORDER BY prod_price DESC, prod_name;


-- Use ORDER BY and LIMIT to find extreme values 
SELECT prod_price 
FROM products 
ORDER BY prod_price DESC LIMIT 1;

\end{lstlisting}


\section{Filter Data}
\subsection{Using the \textsf{WHERE} clause}
Supported operators by \textsf{WHERE}:
\begin{table}
	\centering
	\begin{tabular}{ |c|c|} 
		\hline
		Operator & Description\\ 
		$=$ & Equality \\ 
		$<>$ & Nonequality \\ 
		$!=$ & Nonequality \\ 
		$<$ &  \\ 
		$<=$ &  \\ 
		$>$ &  \\ 
		$>=$ &  \\ 
		$BETWEEN$ & Between two specified values \\ 
		\hline
	\end{tabular}
    \caption{\textsf{WHERE} operators}
    \end{table}
\FloatBarrier


\begin{lstlisting}
-- Check against one single value 
SELECT prod_name, prod_price 
FROM products 
WHERE prod_name = 'fuses';

-- Check against a range of values 
SELECT prod_name, prod_price 
FROM products 
WHERE prod_price BETWEEN 5 AND 10;


-- Check for NO value
SELECT cust_id 
FROM custormer 
WHERE cust_email IS NULL;
\end{lstlisting}

\subsection{Using \textsf{WHERE} with \textsf{AND, OR, IN, NOT IN}}

\begin{lstlisting}
-- and 
SELECT prod_id, prod_price, prod_name
FROM products 
WHERE prod_id = 1000 AND prod_price < 10;

-- or
SELECT prod_id, prod_price, prod_name
FROM products 
WHERE prod_id = 1000 OR prod_price < 10;

-- use parentheses to group if both 'and' and 'or' are used 
SELECT prod_id, prod_price, prod_name
FROM products 
WHERE (prod_id = 1000 or prod_id = 1001) AND prod_price < 10;


-- IN , NOT IN 
SELECT prod_id, prod_price, prod_name
FROM products 
WHERE prod_id IN (1000, 1001)
ORDER BY prod_name;
\end{lstlisting}
\subsection{Wildcard Filtering}

\begin{enumerate}
\item \textsf{LIKE} operator must be used to instruct SQL that the following serach pattern is to be compared using a wildcard match rather than a straight equality match. 

\item ``\%'' matches any number(0, 1 or more) of occurrences of any character. 

\item ``\_'' matches ONE single character. 
\end{enumerate}
\begin{lstlisting}
-- retrieve any value that starts with jet 
SELECT prod_id, prod_price, prod_name
FROM products 
WHERE prod_name LIKE 'jet%'; 

-- retrieve any value that contains jet 
SELECT prod_id, prod_price, prod_name
FROM products 
WHERE prod_name LIKE '%jet%'; 


-- retrieve any value that starts with s and ends with e 
SELECT prod_id, prod_price, prod_name
FROM products 
WHERE prod_name LIKE 's%e'; 


-- retrieve any value that has the second to fourth character as 'jet' 
SELECT prod_id, prod_price, prod_name
FROM products 
WHERE prod_name LIKE '_jet'; 

\end{lstlisting} 

\section{Searching Using Regular Expressions}
\subsection{Basic Character Matching}

\textsf{REGEXP} can be used to match within column values, while \textsf{LIKE} will match the entire column, i.e., if '1000' is only part of the column value, \textsf{LIKE} wouldn't return it while \textsf{REGEXP} will. 

\begin{lstlisting}
-- Match text '1000'
SELECT prod_name 
FROM products 
WHERE prod_name REGEXP '1000'
ORDER BY prod_name;

-- Use '.' to match any single character 
SELECT prod_name 
FROM products 
WHERE prod_name REGEXP '.000'
ORDER BY prod_name;

-- Perform the 'or' matches by '|'
SELECT prod_name 
FROM products 
WHERE prod_name REGEXP '1000|2000|3000'
ORDER BY prod_name;


-- Match one of several specific characters by '[]'
SELECT prod_name 
FROM products 
WHERE prod_name REGEXP '[123] Ton'
ORDER BY prod_name;
-- pattern 1 Ton, 2 Ton or 3 Ton will be matched, and it's a shorthand of [1|2|3]


-- '[]' is needed to tell the or is on 1 2 3 
SELECT prod_name 
FROM products 
WHERE prod_name REGEXP '1|2|3 Ton'
ORDER BY prod_name;
-- means '1' or '2' or '3 Ton' 

-- Negate the matches of several characters  by '^'

SELECT prod_name 
FROM products 
WHERE prod_name REGEXP '[^123] Ton'
ORDER BY prod_name;

\end{lstlisting}


\subsection{Match Ranges and Special Characters}
\begin{enumerate}
\item '-' can be used to define a range. For example, [1-9], [a-z].

\item To match special characters such as '-' used above, need to be preceded by '\verb|\\|'. For example, need to use '\verb|\\|.' to search for '.'. This process to add \verb|\\| is called \emph{escaping}.
\end{enumerate}

\begin{lstlisting}
-- Match range.
SELECT prod_name 
FROM products 
WHERE prod_name REGEXP '[1-5] Ton'
ORDER BY prod_name;


-- Match special character. 
SELECT vend_name 
FROM vendors 
WHERE vend_name REGEXP '\\.'
ORDER BY vend_name;
\end{lstlisting}
\FloatBarrier

\subsection{Matching Character Classes or Multiple Instances}

\begin{table}[h]
	\centering
	\begin{tabular}{ |c|c|} 
		\hline
		Class & Description\\ 
		$[:alnum:]$ & Any letter or digit, [a-zA-Z0-9] \\ 
		$[:alpha:]$ & [a-zA-Z] \\ 
		$[:blank:]$ & Space or Tab \verb|[\\t]| \\ 
		$[:digit:]$ & [0-9] \\ 
		\vdots & \vdots\\
		\hline
	\end{tabular}
	\caption{Character Classes}
\end{table}
\FloatBarrier


\begin{table}[h]
	\centering
	\begin{tabular}{ |c|c|} 
		\hline
		Metacharacter & Description\\ 
		$*$ & 0 or more matches \\ 
		$+$ & 1 or more  matches\\ 
		$?$ & 0 or 1 matches \\ 
		${n}$ & n matches \\ 
		${n, }$ & No less than n matches \\ 
		${n, m}$ & n - m matches \\ 
		\hline
	\end{tabular}
	\caption{Repetition Meta-characters}
\end{table}

\begin{lstlisting}
-- '\\(' matches '(', 'sticks?' matches both 'stick' and 'sticks' as '?' matches 0 or 1 occurrence of WHATEVER IT FOLLOWS.
SELECT prod_name 
FROM products 
WHERE prod_name REGEXP '\\([0-9] sticks?\\)'
ORDER BY prod_name;

-- Match four consecutive digits 
SELECT prod_name 
FROM products 
WHERE prod_name REGEXP '[[:digit:]]{4}'
ORDER BY prod_name;


-- Or equivalently,
SELECT prod_name 
FROM products 
WHERE prod_name REGEXP '[0-9][0-9][0-9][0-9]'
ORDER BY prod_name;
\end{lstlisting}

\subsection{Anchors}

\begin{table}[h]
	\centering
	\begin{tabular}{ |c|c|} 
		\hline
		Metacharacter & Description\\ 
		$\verb|^|$ & Start of TEXT \\ 
		 \$& End of TEXT\\ 
		$[[:<:]]$ & Start of WORD\\ 
		$[[:>:]]$ & End of WORD \\ 
		\hline
	\end{tabular}
	\caption{Anchors}
\end{table}


\begin{lstlisting}

-- Match start of a . or any digit
SELECT prod_name 
FROM products 
WHERE prod_name 'REGEXP ^[0-9\\.]'
ORDER BY prod_name;
\end{lstlisting}


\section{Data Manipulation and Summarization}

\subsection{Creating Calculated Fields}

\subsubsection{Concatenate Fields} 
\textsf{Concat} function is used to form a single long value by joining values together. For most DBMS, however, '+' or '$||$' are used for concatenation.
\begin{lstlisting}

-- Concatenate name and country to be a single column
SELECT Concat(vend_name, '(', vend_country, ')')
FROM vendors
ORDER BY vend_name;
\end{lstlisting}

\subsection{Using Aliases}
\textsf{AS} is used. 
\begin{lstlisting}
-- Assign a name 
SELECT Concat(vend_name, '(', vend_country, ')') AS vend_title
FROM vendors
ORDER BY vend_name;


-- Mathmatical calculations such as +, -, *, / are allowed when retriving the data.
\end{lstlisting}

\subsection{Data Manipulation}
\begin{table}[h]
	\centering
	\begin{tabular}{ |c|c|} 
		\hline
		\textbf{Text Manipulation} & Description\\ 
		\hline
		Left() & Returns characters from left of string \\ 
		Length()& Returns the length of the string\\ 
		Locate() & Finds a substring within a string\\ 
		Lower() & Coverts string to lowercase \\ 
		Upper() & Converts string to uppercase\\
		LTrim() & Trims white space from left of string\\
		RTrim() & Trims white space from right of string\\
		Soundex() & Returns a string's soundex values\\
		\hline
		\textbf{Numeric Manipulation Functions} & Description\\
		\hline
		Abs() & Absolute value \\
		Cos() & Cos(x)\\
		Exp() & Exponential value \\
		Mod() & Remainder of a division operation \\
		Pi() & pi\\
		Rand() & Random number \\
		Sin() & sin()\\
		Sqrt() & Square root \\
		Tan() & tan()\\
		\hline
	\end{tabular}
	\caption{Manipulation Functions}
\end{table}
\FloatBarrier

\begin{lstlisting}
SELECT vend_name, UPPER(vend_name) AS vend_name_upper
FROM vendors 
ORDER BY vend_name;

-- Use Soundex to return values which sound similar with Y Lie, for example, Y. Lie
SELECT vend_name
FROM vendors 
WHERE Soundex(vend_name) = Soundex('Y Lie)
ORDER BY vend_name;

-- It's a good practice to use Date() for dates to avoid confusion, and date needs to have the format yyyy-mm-dd
SELECT vend_name
FROM vendors 
WHERE Date(order_date) = '2005-09-01'
ORDER BY vend_name;

-- Method 1 to retrieve all orders in one certain month 
SELECT cust_id, order_num
FROM orders 
WHERE Date(order_date) BETWEEN '2005-09-01' AND '2005-09-30';


-- Method 2
SELECT cust_id, order_num
FROM orders 
WHERE Year(order_date) = 2005 and Month(order_date) = 9;
-- NO QUOTES!
\end{lstlisting}


\subsection{Summarization}
\begin{table}[h]
	\centering
	\begin{tabular}{ |c|c|} 
		\hline
		Aggregate Functions & Description\\ 
		\hline
		AVG() & Average \\ 
		COUNT() & Number of specified rows\\ 
		MAX() & Max\\ 
		MIN() & Min\\ 
		SUM() & Sum\\ 
		\hline
	\end{tabular}
	\caption{Aggregation Functions}
\end{table}


\begin{lstlisting}
-- Select all rows 
SELECT COUNT(*) AS num_cust
FROM custormers;


-- Count specified rows 
SELECT COUNT(cust_email) AS num_cust
FROM custormers;
\end{lstlisting}

\subsection{Summarization on Grouped Data}

\begin{table}[h]
	\centering
	\begin{tabular}{ |c|c|} 
		\hline
		Clause & Description\\ 
		\hline
		GROUP BY & Group specification \\ 
		HAVING & Group level filtering\\ 
		WHERE & Row level filtering\\ 
		\hline
	\end{tabular}
	\caption{Clause}
\end{table}
\begin{lstlisting}
-- Calculate the COUNT by vend_id 
SELECT vend_id, COUNT(*) AS num_prods 
FROM products
GROUP BY vend_id;


-- Filter groups with HAVING
SELECT cust_id, COUNT(*) AS orders 
FROM orders 
GROUP BY cust_id 
HAVING COUNT(*) >= 2;


-- Note the ORDER BY posistion 
SELECT cust_id, COUNT(*) AS orders 
FROM orders 
GROUP BY cust_id 
HAVING COUNT(*) >= 2
ORDER BY cust_id;
\end{lstlisting}


\section{Subqueries and Joining Tables}

\subsection{Subquery}
Subquery is query that is embedded into other queries. Subqueries are always processed starting with the INNERMOST \textsf{SELECT} statement and working outward. 

Subqueries can also be done using JOINS in next Chapter. 

It can be used to do: 
\begin{itemize}
\item filtering

\item calculated fields 
\end{itemize}

\begin{lstlisting}
-- Filtering, indentation can be used for easier understanding of the code. 
SELECT cust_id 
FROM orders 
WHERE order_num IN (
                    SELECT order_num 
                    FROM orderitems
                    WHERE prod_id = 'TNT2');
                    
-- Calculated field 
SELECT cust_name, 
       cust_state,
       (SELECT COUNT(*) 
       FROM orders
       WHERE orders.cust_id = customers.cust_id) AS orders
FROM customers
ORDER BY cust_name;                    
\end{lstlisting}

TIPS: When need to use subqueries, it's safe to do so incrementally in much the same way as SQL processes them. 


\subsection{Join Tables}
Foreign key is used to join tables. 

\begin{itemize}
\item Join can be created by simply specifying all tables and how they are related using \textsf{WHERE}. 

\item ``INNER JOIN'' and  ``ON'' cane be used together to do the job. 

\item Joining multiple tables can be done by first list all the tables and then define the relationship. 

\item OUTER JOIN: joins including rows that have no associated rows in the related table. ``LEFT OUTER JOIN'' , RIGHT OUTER JOIN'' can be used. 

\item JOINs can be combined with aggregate functions. 
\end{itemize}
\begin{lstlisting}

-- Create a join by WHERE 
SELECT vend_name, prod_name 
FROM  vendors, products 
WHERE vendors.vend_id = products.vend_id
ORDER BY vend_name, prod_name;


-- Use ``INNER JOIN'' with ``ON'' to do the same as above 
SELECT vend_name, prod_name 
FROM  vendors INNER JOIN products 
ON vendors.vend_id = products.vend_id
ORDER BY vend_name, prod_name;


-- Join multiple tables 
SELECT cust_name, cust_contact
FROM customers, orders, orderitems
WHERE customers.cust_id = orders.cust_id 
AND orderitems.order_num = orders.order_num
AND orderitems.prod_id = 'TNT2';


-- OUTER JOIN 
SELECT customers.cust_id, orders.order_num
FROM customers LEFT OUTER JOIN orders 
ON customers.cust_id = orders.cust_id;
-- Retrieve all customers, even no orders(order_num is NULL) placed.
-- LEFT OUTER JOIN selects all rows on the left. 


-- Aggregate function , inner join 
SELECT customers.cust_name,
       customers.cust_id,
       COUNT(orders.order_num) AS num_ord
FROM customers INNER JOIN orders
ON customers.cust_id = orders.cust_id
GROUP BY customers.cust_id;
       
-- Outer join        
SELECT customers.cust_name,
customers.cust_id,
COUNT(orders.order_num) AS num_ord
FROM customers LEFT OUTER JOIN orders
ON customers.cust_id = orders.cust_id
GROUP BY customers.cust_id;

\end{lstlisting}


\subsection{Combine Queries}
Performing multiple queries and return the results as a single query result set is know as \emph{unions} or \emph{compound queries}. 
\\~\\
\textbf{Union Rules}: 
\begin{enumerate}
	
\item A UNION must be comprised of two or more SELECT statements, each separated by keyword UNION.

\item Each query in a UNION must contain the same columns, expressions, or aggregate functions. 

\item Column datatypes must be compatible.  
\end{enumerate}

\textsf{UNION} can eliminate duplicated rows by default. If one wants all occurrences of all matches returned, use \textsf{UNION ALL}. 

\begin{lstlisting}
--  Use of UNION
SELECT vend_id, prod_id, prod_price
FROM products
WHERE prod_price <= 5
UNION
SELETCT vend_id, prod_id, prod_price
FROM products
WHERE vend_id IN (1001, 1002);

-- Or equivalently,
SELECT vend_id, prod_id, prod_price
FROM products
WHERE prod_price <= 5 
OR vend_id IN (1001, 1002);

-- Some duplicated rows will return (if there are any).
SELECT vend_id, prod_id, prod_price
FROM products
WHERE prod_price <= 5
UNION ALL
SELETCT vend_id, prod_id, prod_price
FROM products
WHERE vend_id IN (1001, 1002);


-- Sorted by ORDER BY at the end of the final SELECT
SELECT vend_id, prod_id, prod_price
FROM products
WHERE prod_price <= 5
UNION
SELETCT vend_id, prod_id, prod_price
FROM products
WHERE vend_id IN (1001, 1002)
ORDER BY vend_id, prod_price;
\end{lstlisting}

\section{Full-Text Searching}

\begin{enumerate}
\item In order to perform full-text searching, the columns to be searched must be indexed and constantly re-indexed as data changes. 

\item \textsf{SELECT} can be used with \textsf{MATCH()} and \textsf{Against()} for searches. 

\item ``WITH QUERY EXPANSION'' can be used within ``Against'' to widen the range of returned full text results. The match is not only based on the original criteria, but also useful words from the selected matched rows. 

\item Boolean mode can be used to specify 1).words to be matched/excluded; 2).ranking hints on which words are more important; 3).etc.
\end{enumerate}

\begin{table}[h]
	\centering
	\begin{tabular}{ |c|l|} 
		\hline
		Privilege & Description\\ 
		\hline
		$+$ & Word must be included. \\ 
		\hline
		$-$ & Word must be excluded. \\ 
		\hline
		$>$ & Word must be included and increase ranking value. \\
		\hline
		$<$ & Word must be included and decrease ranking value. \\
		\hline
		$()$ & Group words into subexpressions. \\  
		\hline
		$\sim$ & Negate a word's ranking value. \\  
		\hline
		$*$ & Wildcard at end of word. \\ 
		\hline
		$" "$ & Define a phrase instead of individual words. \\ 

		\hline
	\end{tabular}
	\caption{Boolean Operators}
\end{table}


\begin{lstlisting}
-- Index the column by FULLTEXT()
CREATE TABLE productnotes 
(
note_id        int        NOT NULL AUTO_INCREMENT,
prod_id        char(10)   NOT NULL,
note_date      datetime   NOT NULL,
note_text      text       NULL,
PRIMARY KEY(note_id),
FULLTEXT(note_text)
) ENGINE = MyISAM;


-- full text search 
SELECT note_text 
FROM productnotes
WHERE Match(note_text) Against('rabbit');

-- or equivalently
SELECT note_text 
FROM productnotes
WHERE note_text LIKE '%rabbit%';


-- Query expansion 
SELECT note_text 
FROM productnotes
WHERE Match(note_text) Against('rabbit' WITH QUERY EXPANSION);
-- More results with the selected useful words will be returned.


-- Boolean mode 
SELECT note_text 
FROM productnotes
WHERE Match(note_text) Against('heavy -rope*' IN BOOLEAN MODE);
-- 'heavy' is matched, and excludeany result that contains rope*(any word beginning with 'rope')

-- use of + 
SELECT note_text 
FROM productnotes
WHERE Match(note_text) Against('+heavy +rope*' IN BOOLEAN MODE);


-- Single quote
SELECT note_text 
FROM productnotes
WHERE Match(note_text) Against('heavy rope' IN BOOLEAN MODE);
-- Match results that contain at least one of heavy and rope.


-- Double quotes 
SELECT note_text 
FROM productnotes
WHERE Match(note_text) Against('"heavy rope"' IN BOOLEAN MODE);
-- Match phrase heavy rope instead of two words. 

--
SELECT note_text 
FROM productnotes
WHERE Match(note_text) Against('+safe +(<combination)' IN BOOLEAN MODE);
-- Match safe and combination, lower the ranking of the latter. 
\end{lstlisting}

\section{Inserting/Updating/Deleting Data}
\subsection{Inserting Data}

\begin{enumerate}
\item \textsf{INSERT INTO\ldots VALUES} syntax is used. 

\item Multiple rows could be inserted at one time, while different \textsf{INSERT} statements separated by semicolon. Also, if column names are identical for different \textsf{INSERT} statements, one need to specify only one time of the column list. 

\item Replace the \textsf{VALUES} part with the normal \textsf{SELECT} statements to insert retrieved data.  
\end{enumerate}


\begin{lstlisting}
-- Insert a single row 
INSERT INTO customers(
  cust_name,
  cust_address,
  cust_city,
  cust_state,
  cust_email)
VALUES(
  'Pep LePew',
  '100 Main Street',
  'Los Angles',
  'CA',
  NULL);
-- First specify the column list to be safe, then specify the values to be inserted, NULL allowed given that the column allows NULL or default value is assigned if NULL.

-- Insert multiple rows, column lists need not to be identical 
INSERT INTO customers(
cust_name,
cust_address,
cust_city,
cust_state,
cust_email)
VALUES(
'Pep LePew',
'100 Main Street',
'Los Angles',
'CA',
NULL);
INSERT INTO customers(
cust_name,
cust_address,
cust_city,
cust_state,
cust_email)
VALUES(
'Pep LePew',
'100 Main Street',
'Los Angles',
'CA',
NULL);



-- If column lists are identical
INSERT INTO customers(
cust_name,
cust_address,
cust_city,
cust_state,
cust_email)
VALUES(
'Pep LePew',
'100 Main Street',
'Los Angles',
'CA',
NULL),
(
'Pep LePew',
'100 Main Street',
'Los Angles',
'CA',
NULL);


-- Insert retrieved data,  
INSERT INTO customers(
cust_name,
cust_address,
cust_city,
cust_state,
cust_email)
SELECT cust_name,
       cust_address,
       cust_city,
       cust_state,
       cust_email
FROM custnew;

\end{lstlisting} 


\subsection{Updating and Deleting Tables}

\begin{enumerate}
\item \textsf{UPDATE \ldots SET} syntax is used. Multiple columns can be updated together with a single \textsf{SET} and column = value pair is separated by a comma. 

\item Deleting a certain column can be accomplished by setting the specific column value to NULL. 

\item \textsf{DELETE FROM} syntax is used to delete a entire row, not a specific column. 

\item \textsf{TRUNCATE TABLE} is used to delete all rows from a table. It drops and recreates the table. 
\end{enumerate}


\begin{lstlisting}

-- Update one column 
UPDATE customers
SET cust_email = 'sql@test.com'
WHERE cust_id = 1000;

-- Update multiple columns
UPDATE customers
SET cust_name = 'SQL',
    cust_email = 'sql@test.com'
WHERE cust_id = 1000;


-- Delete a column value 
UPDATE customers
SET cust_email = NULL
WHERE cust_id = 1000;

-- Delete a row 
DELETE FROM customers 
WHERE cust_id = 10000;


-- Delte the entire rows 
TRUNCATE TABLE customers;  
\end{lstlisting}


\section{Creating and Manipulating Tables}

\subsection{Create Table}
\begin{enumerate}
\item \textsf{CREATE TABLE tablename \ldots} is used to create a table, and name and definition of different columns is separated by a comma. 

\item NULL states of a column is specified at the table definition by ``NULL'' or ``NOT NULL''. 

\item PRIMARY KEY must be unique and defined when creating the table.

\item \textsf{AUTO\_INCREMENT} is used to tell SQL that the column is automatically incremented each time a row is added. Only one column is allowed to be auto-incremented, and it must be indexed(for example, making it a primary key).

\item \textsf{DEFAULT} is used to specify the default values of a column in \textsf{CREATE TABLE}. 

\item \textsf{ENGINE = } is used to specify the internal engine that SQL uses to manage and manipulate data. 
\end{enumerate}

\begin{lstlisting}
-- Create a table 
CREATE TABLE vendors 
(
vend_id     int       NOT NULL AUTO_INCREMENT,
vend_name   char(50)  NOT NULL DEFAULT test,
vend_city   char(50)  NULL, 
PRIMARY KEY(vend_id, vend_name)
) ENGINE = InnoDB;
\end{lstlisting}

\subsection{Manipulate Table}

\begin{itemize}
\item Update table definition by \textsf{ALTER TABLE tablename ADD/DROP \ldots}. 

\item Delete table by \textsf{DROP TABLE tablename}.

\item Rename table by \textsf{RENAME TABLE tablename1 TO tablename2}.
\end{itemize}

\begin{lstlisting}
-- Update tables : add 
ALTER TABLE vendors 
ADD vend_phone CHAR(20);


-- Update tables: drop 
ALTER TABLE vendors 
DROP COLUMN vend_phone;

-- Define foreign keys 
ALTER TABLE orderitems 
ADD CONSTRAINT fk_orderitems_orders
FOREIGN KEY (order_num) REFERENCES orders (order_num);

-- Drop table 
DROP TABLE customers;

-- Rename 
RENAME TABLE customers1 TO customers2,
             customers3 TO customers4;

\end{lstlisting}
\section{Using Views and Stored Procedures}
\subsection{Views}
Views are virtual tables. It contains queries that dynamically retrieve data when used. 

The common uses are: 
\begin{enumerate}
\item Reuse SQL statement. 
\item Simplify complex SQL operations. 
\item To expose part of a table instead of complete tables. 
\item To secure data.
\item Return data formatted and presented differently from their underlying tables. 
\end{enumerate}

Views are created using \textsf{CREATE VIEW viewname}, and dropped by \textsf{DROP VIEW viewname}.

\begin{lstlisting}
-- Simplify complex joins : step 1: create a view 
CREATE VIEW productcustomers AS 
SELECT cust_name, cust_contact, prod_id
FROM custormers, orders, orderitems
WHERE customers.cust_id = orders.cust_id
AND   orderitems.order_num = orders.order_num;

-- step2 : retrieve data using view 
SELECT cust_name, cust_contact
FROM productcustomers
WHERE prod_id = 'TNT2';

-- Reformat data 
CREATE VIEW vendorlocations AS
SELECT Concat(RTrim(vend_name), '(', RTrim(vend_country), ')')
AS vend_title
FROM vendors
ORDER BY vend_name;


SELECT * 
FROM vendorlocations;


-- Filter unwanted data , for example, filter out data without email address
CREATE VIEW customeremaillist AS 
SELECT cust_id, cust_name, cust_email
FROM customers
WHERE cust_email IS NOT NULL;

SELECT * 
FROM customeremaillist;


-- Use views with calculated fields 
CREATE VIEW orderitemexpand AS 
SELECT order_num, 
       prod_id, 
       quantity, 
       item_price,
       quantity*item_price AS expanded_price
FROM orderitems;

SELECT * 
FROM orderitemsexpand
WHERE order_num = 20005; 
\end{lstlisting}


\subsection{Stored Procedures}
Stored procedures are collections of one or more SQL statements saved for future use. 

\begin{enumerate}
\item Stored procedures are executed by \textsf{CALL}. It will not display data by only calling it, need to use \textsf{SELECT} to retrieve data. 

\item Stored procedures are created by \textsf{CREATE PROCEDURE \ldots BEGIN \ldots END}.


\item Default SQL delimiter is \verb|;| while the command-line utility also uses it as a delimiter. One can first change the command utility delimiter to something else and change it back later. For example, \verb|DELIMITER // | \ldots \verb| DELIMITER ; |. Any character expcept for \verb|\| can be a delimiter in SQL.

\item Procedure can have parameters. 
\end{enumerate}


\begin{lstlisting}
-- Create a procedure 
CREATE PROCEDURE productpricing()
BEGIN 
     SELECT Ave(prod_price) AS priceaverage
     FROM products;
END;

-- Call the procedure 
CALL productpricing();
-- the trailing () is required. 


-- Drop 
DROP PROCEDURE productpricing; 
-- No trailing. 

-- With parameters 
CREATE PROCEDURE productpricing(
        OUT pl DECIMAL(8,2),
        OUT ph DECIMAL(8,2),
        OUT pa DECIMAL(8,2)
        )
      BEGIN
         SELECT Min(prod_price) INTO pl
         FROM products;
         SELECT Max(prod_price) INTO ph
         FROM products;
         SELECT Avg(prod_price) INTO pa
         FROM products;
       END;
-- OUT: specify that this parameter is used to send a value out of stored procedure. Can also use IN to pass values to stored procedure; INOUT to pass parameters to and from stored procedures.


-- call the procedure with parameters 
CALL productpricing(
                    @pricelow,
                    @pricehigh.
                    @priceaverage);      

-- Retrieve data 
SELECT @priceaverage, @pricelow;

-- Use both IN and OUT 
CREATE PROCEDURE ordertotal(
       IN onumber INT 
       OUT ototal DECIMAL(8,2)
       )
    BEGIN
       SELECT SUM(item_price*quantity) 
       FROM orderitems
       WHERE order_num = onumber INTO ototal;
    END;
    
-- Call the procedure 
CALL ordertotal(20005, @total);
SELECT @total;


-- Example, add tax for total only for specific customers, business rules and intelligent processing will be included. 

-- Name: ordertotal
-- Parameters: onumber = order number
--             taxable = 1 taxable, 0 not 
--             ototal = order total 

CREATE PROCEDURE ordertotal(
       IN onumber INT 
       IN taxable BOOLEAN 
       OUT ototal DECIMAL(8,2)
       ) COMMENT 'Obtain order total with optional tax'
BEGIN
-- Declare variable for total ,local variables
DECLARE total DECIMAL(8,2);
-- Tax percentage 
DECLARE taxrate INT DEFAULT 6;

-- Get order total 
SELECT SUM(item_price*quantity)
FROM orderitems
WHERE order_num = onumber
INTO total; 

IF taxable THEN SELECT total + (total * taxrate /100) 
INTO total;
END IF;

SEELECT total INTO ototal;
END;


-- Call 
CALL ordertotal(20005, 1, @total);
SELECT @total;

-- Inspect stored procedures: display the statement used to create a procedure 
SHOW CREATE PROCEDURE ordertotal;

-- Inspect: a list of ALL stored procedures including details 
SHOW PROCEDURE STATUS  'ordertotal';

-- Show only specified procedures 
SHOW PROCEDURE STATUS LIKE 'ordertotal';

\end{lstlisting}
\section{Cursors and Triggers}
\subsection{Cursors}
A cursor is a database query stored on the SQL server. It's not a SELECT statement, but the result set retrieved by that statement.

\begin{enumerate}
\item A cursor must be declared/defined before it's used.
\item The cursor must be opened for use. This process retrieves data. 
\item With cursor populated with data, individual rows can be fetched/retrieved. 
\item The cursor must be closed when it's done. 
\end{enumerate} 


\begin{lstlisting}
CREATE PROCEDURE processorders()
BEGIN 
-- Declare local variables 
DECLARE done BOOLEAN DEFAULT 0;
DECLARE o INT; 
DECLARE t DECIMAL(8,2);

-- Declare the cursor 
DECLARE ordernumbers CURSOR 
FOR 
SELECT order_num FROM orders;

-- Declare continue handler 
DECLARE CONTINUE HANDLER FOR SQLSTATE '02000' SET done = 1;


-- Create table to store results if not exists 
CREATE TABLE IF NOT EXISTS ordertotals
(order_num INT, total DECIMAL(8,2));

-- Open cursor 
OPEN ordernumbers;

-- Loop through all rows 
REPEAT

-- Get order numbers 
FETCH ordernumbers INTO o;

-- Get the total for this order 
CALL ordertotal(o, 1, t);

-- Insert order number and order total into the table 
INSERT INTO ordertotals(order_num, total) 
VALUES(o, t);

-- End of loop 
UNTIL done END REPEAT;

-- Close the cursor 
CLOSE ordernumbers;

END:



-- Display the results in the table 
SELECT * 
FROM ordertotals;

\end{lstlisting}
\begin{itemize}
\item CONTINUE HANDLER will be executed when a condition occurs, while SQLSTATE '02000' is a not found condition, so it occurs when REPEAT cannot continue(no more rows to loop). 

\item A table \emph{ordertotals} is created to fetch each order number. 
\end{itemize}


\subsection{Triggers}
A trigger is a statement(or a group of statements enclosed with BEGIN and END statements) that are automatically executed by SQL in response to any of these statements: 
\begin{itemize}
\item DELETE: within DELETE trigger code, can refer to a virtual table named ``OLD'' to access the rows being deleted. 

\item INSERT: within INSERT trigger code, can refer to a virtual table named ``NEW'' to access rows being inserted. 

\item UPDATE: within UPDATE trigger code, can refer to a virtual table named ``OLD'' to access the previous values and ``NEW'' to access the new updated values. 

\end{itemize}

Four piece of information is needed to create a trigger: 
\begin{enumerate}
\item Unique trigger name.
\item The table to which the trigger is to be associated.
\item The action that the trigger should respond to(the three mentioned above).
\item When the trigger should be executed. 
\end{enumerate}

\begin{enumerate}
\item \textsf{CREATE TRIGGER triggername AFTER/BEFORE \ldots FOR EACH ROW \ldots;} syntax is used. 

\item Drop the trigger by \textsf{DROP TRIGGER triggername;} Triggers can not be updated or overwritten. 
\end{enumerate}



\begin{lstlisting}

-- Create trigger 
CREATE TRIGGER newproduct AFTER INSERT ON products 
FOR EACH ROW SELECT 'Product added';
-- Trigger is executed after a successful INSERT, then for each added row, display "Product added". 

-- Drop trigger
DROP TRIGGER newproduct;


-- Use NEW to access the order number being inserted 
CREATE TRIGGER neworder AFTER INSERT ON orders 
FOR EACH ROW SELECT NEW.order_num;
-- Return the new order number after inserted 


-- Use OLD to save rows to be deleted into an active table 
CREATE TRIGGER deleteorder BEFORE DELETE ON orders 
FOR EACH ROW 
BEGIN 
    INSERT INTO archive_orders(order_num, order_date, cust_id)
    VALUES(OLD.order_num, OLD.order_date, OLD.cust_id);
END;
-- Use BEFORE DELTE has advantage that : if for some reason the order could not be archived, the DELETE will be aborted.




-- Use UPDATE 
CREATE TRIGGER updatevendor BEFORE UPDATE ON vendors
FOR EACH ROW SET NEW.vend_state = UPPER(NEW.vend_state);
\end{lstlisting}

This is added for test.
A second change
A third change
\end{document}
